\documentclass[11pt,a4paper,DIV=10,parskip=half,BCOR=0mm]{scrartcl}
\usepackage[utf8]{inputenc}
\usepackage[T1]{fontenc}
\usepackage[ngerman]{babel}
\usepackage{scrpage2} % pagestyle
\usepackage{geometry} % Seitenränder
\usepackage{amsmath}
\usepackage{amsfonts}
\usepackage{amssymb}
\usepackage{amsthm}
\usepackage{multicol}
\usepackage{graphicx}
\usepackage{float}
\usepackage{marginnote} % Randnotizen
\usepackage{tabularx} % Tabelle mit spaltentyp 'X' für feste breite
\columnsep6mm 
% Seitenränder einstellen: marginparwidth wichtig für Randnotizen!
\geometry{a4paper, top=20mm, left=35mm, right=50mm, bottom=30mm,
	headsep=10mm, footskip=12mm, marginparwidth=45mm, marginparsep=3mm}
	
% Fußzeile
\pagestyle{scrheadings}
\ifoot{ \hrule \vspace{1mm} \small{Erstellt am: \today \ von C. van Heteren-Frese \\
%Grafik: CC BY-SA 3.0 Fiona Macintosh
}} 
\begin{document}
\pagenumbering{gobble}
% === HEADER ====
\newcolumntype{T}{>{\Large\centering\arraybackslash}X}
\setlength{\tabcolsep}{3mm} % kein Innenrand bei Spaleten
\noindent
\begin{tabularx}{\textwidth}{|l|T|l|}
\hline
Datum: & \rule{0pt}{5mm} \textbf{\textsf{Thema: Interessengruppen des Smart Grid}} & Blatt-Nr. \\
&&\\
\hline
\end{tabularx}
% === HEADER END ===

%
% === HEADER ====
\setlength{\tabcolsep}{0mm} % kein Innenrand bei Spaleten
\begin{tabularx}{\linewidth}{lXr}
{\Large\textsf{\textbf{Gruppe A:} Firma \textit{Alleswisser GmbH \& Co. KG}}} & &
\includegraphics[scale=0.25]{images/euro}\\ % Logo rechtsbündig
\hline
\end{tabularx}
% === HEADER END ===
%
\subsection*{Hintergrund}
Du bist Manager der privaten Firma \textit{Alleswisser GmbH \& Co. KG}.
Die Firma von dir und deinen Kollegen ermittelt \marginnote{\footnotesize \textbf{personenbezogene Daten} können einer bestimmten Person zugeordnet werden} personenbezogene Daten und filtert daraus spezielle Informationen (z.B. wann und wie lange bei
Privatpersonen zu Hause das Licht eingeschaltet ist). Aus diesen
Informationen erstellt ihr ein spezielles Kundenprofil \marginnote{\footnotesize \textbf{Kundenprofil} enthält Eigenschaften und Verhaltensweisen}und verkauft es
anschließend gezielt an andere Firmen.

\textbf{Beispiel: }Eine Firma verkauft einen speziellen, wach machenden Energy-
Tee „Immerwach“. Diese Firma möchte bei euch ganz spezielle Adressen
einkaufen: und zwar von Leuten, die nachts lange wach sind! An diese
Adressen soll gezielt Werbung für den Tee geschickt werden.
\subsection*{Infos}
\begin{enumerate}
	\item[•]Die Daten, die ihr verarbeitet, bezieht ihr -- wenn möglich -- von den Netzbetreibern (z.B.
	Deutsche Telekom)
	\item[•]Falls ihr Daten braucht, die nicht ganz legal zu bekommen sind, holt
	ihr euch diese auch schon mal von Informatikstudenten oder
	anderen Zwischenhändlern. Woher diese die Daten haben,
	interessiert euch nicht.
\end{enumerate}
\subsection*{Eure Position}
\begin{enumerate}
	\item[•]Ihr seid \textbf{für} das \textit{Smart Grid}
\end{enumerate}
\subsection*{Eure Gründe}
\begin{enumerate}
	\item[•]Euer ganzes Geschäft beruht auf dem Handel mit Daten. Gibt es
	kein \textit{Smart Grid}, seid ihr pleite! Das müsst ihr in der Talk Show
	verhindern!
	\item[•] Euer Hauptargument: \glqq Wenn man nichts zu verbergen hat, muss man sich keine Gedanken um den Datenschutz machen!\grqq
\end{enumerate}
\end{document}